% =====================================================================
\section{Steepest slope of the Lorentzian line}
\label{app:fisher}
% =====================================================================

This appendix derives the best information per measurement, Eq.~\eqref{eq:istar}, for the Lorentzian line of
Eq.~\eqref{eq:signal}. Written in terms of the detuning $u=x-\theta$ of the measurement from the resonance, the
line depends on $x$ and $\theta$ only through $u$, so a shift of the resonance acts like the opposite shift of the
measurement, and the slope of Eq.~\eqref{eq:fisher} is
\begin{equation*}
  \frac{\partial\mu}{\partial\theta}=-\frac{\partial\mu}{\partial u}=-\frac{2C\gamma^2u}{(\gamma^2+u^2)^2}.
\end{equation*}
It vanishes at $u=0$ and as $u\to\pm\infty$. Setting $\partial^2\mu/\partial u^2=0$ gives
$(\gamma^2+u^2)^2=4u^2(\gamma^2+u^2)$, i.e.\ $u^2=\gamma^2/3$: the slope is steepest at $|u|=\gamma/\sqrt3$, where
$\gamma^2/(\gamma^2+u^2)=3/4$ and the dip is $3/4$ deep. There
\begin{equation*}
  \begin{aligned}
  \max_x\Big|\frac{\partial\mu}{\partial\theta}\Big|&=\frac{2C\gamma^2\,(\gamma/\sqrt3)}{(4\gamma^2/3)^2}=\frac{9C}{8\sqrt3\,\gamma}\\
  &=\frac{9C}{4\sqrt3\,\Gamma}\approx1.30\,\frac{C}{\Gamma},
  \end{aligned}
\end{equation*}
with $\Gamma=2\gamma$. Inserting this slope into Eq.~\eqref{eq:fisher}, with $\mathrm{SNR}=C/\sigma$,
\begin{equation*}
  I_{\rm max}=\frac{k}{\sigma^2}\Big(\frac{9C}{4\sqrt3\,\Gamma}\Big)^{\!2}=\frac{81}{48}\,\frac{k\,C^2}{\sigma^2\Gamma^2}
  =\frac{27}{16}\,k\,\frac{\mathrm{SNR}^2}{\Gamma^2}.
\end{equation*}
